\documentclass[11pt, a4paper]{article}

\usepackage[utf8]{inputenc}
\usepackage[T1]{fontenc}
\usepackage[french]{babel}
\usepackage{geometry}
\usepackage{tikz}

% TikZ libraries for flowcharts
\usetikzlibrary{shapes.geometric, arrows, calc, positioning}

\geometry{
    hmargin=2.5cm,
    vmargin=2.5cm
}

% Define TikZ styles for flowchart elements
\tikzstyle{startstop} = [rectangle, rounded corners, minimum width=3cm, minimum height=1cm, text centered, draw=black, fill=red!30]
\tikzstyle{io} = [trapezium, trapezium left angle=70, trapezium right angle=110, minimum width=3cm, minimum height=1cm, text centered, draw=black, fill=blue!30]
\tikzstyle{process} = [rectangle, minimum width=3cm, minimum height=1cm, text centered, draw=black, fill=orange!30]
\tikzstyle{decision} = [diamond, minimum width=3cm, minimum height=1cm, text centered, draw=black, fill=green!30, aspect=2]
\tikzstyle{arrow} = [thick,->,>=stealth]

\title{TD1 - Les Organigrammes}
\author{Adel HACHACHE}
\date{3 décembre 2025}

\begin{document}

\maketitle

\section*{Exercice 1: Calcul du Factoriel}

\paragraph{Énoncé}
Créez un organigramme qui calcule le factoriel d'un nombre entier non négatif n. Le factoriel d'un nombre n, noté n!, est le produit de tous les entiers positifs inférieurs ou égaux à n.

\paragraph{Instructions}
\begin{enumerate}
    \item Commencez par initialiser le résultat à 1.
    \item Utilisez une boucle pour multiplier le résultat par chaque nombre de 1 à n.
    \item Affichez le résultat final.
\end{enumerate}

\begin{center}
\begin{tikzpicture}[node distance=1.5cm]

% Nodes
\node (start) [startstop] {Début};
\node (in1) [io, below of=start] {Lire n};
\node (proc1) [process, below of=in1] {fact $\leftarrow$ 1};
\node (proc2) [process, below of=proc1] {i $\leftarrow$ 1};

% This coordinate is the point where the loop will rejoin the main flow
\coordinate (join) at ([yshift=-1cm]proc2.south);

\node (dec1) [decision, below of=join] {i $\le$ n ?};
\node (proc3) [process, right of=dec1, xshift=2.5cm] {fact $\leftarrow$ fact * i};
\node (proc4) [process, below of=proc3] {i $\leftarrow$ i + 1};
\node (out1) [io, below of=dec1, yshift=-1.5cm] {Afficher fact};
\node (stop) [startstop, below of=out1] {Fin};

% Arrows
\draw [arrow] (start) -- (in1);
\draw [arrow] (in1) -- (proc1);
\draw [arrow] (proc1) -- (proc2);
\draw [arrow] (proc2) -- (join);
\draw [arrow] (join) -- (dec1);

\draw [arrow] (dec1.east) -- node[above] {Oui} (proc3);
\draw [arrow] (dec1.south) -- node[right] {Non} (out1);
\draw [arrow] (proc3) -- (proc4);

% The loop back arrow, drawn cleanly to the join point
\draw [arrow] (proc4.east) -- ++(0.5,0) |- (join);

\draw [arrow] (out1) -- (stop);

\end{tikzpicture}
\end{center}

\newpage

\section*{Exercice 2: Détection de Nombres Pairs et Impairs}

\paragraph{Énoncé}
Construisez un organigramme qui détecte si un nombre entier n est pair ou impair.

\paragraph{Instructions}
\begin{enumerate}
    \item L'organigramme doit commencer par la saisie du nombre n.
    \item Utilisez un test conditionnel pour vérifier la parité de n en utilisant le reste de la division par 2.
    \item Affichez si le nombre est pair ou impair.
\end{enumerate}

\begin{center}
\begin{tikzpicture}[node distance=2cm]
\node (start) [startstop] {Début};
\node (in1) [io, below of=start] {Saisir n};
\node (dec1) [decision, below of=in1, yshift=-1cm] {n mod 2 == 0 ?};
\node (out_even) [io, right of=dec1, xshift=3cm] {Afficher "Pair"};
\node (out_odd) [io, left of=dec1, xshift=-3cm] {Afficher "Impair"};
\node (stop) [startstop, below of=dec1, yshift=-2cm] {Fin};

\draw [arrow] (start) -- (in1);
\draw [arrow] (in1) -- (dec1);
\draw [arrow] (dec1.east) -- node[above] {Oui} (out_even);
\draw [arrow] (dec1.west) -- node[above] {Non} (out_odd);
\draw [arrow] (out_odd) -- (stop);
\draw [arrow] (out_even) -- (stop);
\end{tikzpicture}
\end{center}

\section*{Exercice 3 : Trouver le Plus Grand de Trois Nombres}

\paragraph{Énoncé}
Développez un organigramme qui prend en entrée trois nombres a, b, et c, et détermine lequel est le plus grand.

\paragraph{Instructions}
\begin{enumerate}
    \item L'organigramme doit comparer les nombres deux à deux.
    \item Utilisez des structures conditionnelles pour identifier le plus grand nombre.
    \item Affichez le plus grand nombre.
\end{enumerate}

\begin{center}
% Cet organigramme utilise une approche séquentielle pour une meilleure lisibilité.
% 1. On initialise 'max' avec la valeur de 'a'.
% 2. On compare 'b' à 'max'. Si 'b' est plus grand, 'max' prend la valeur de 'b'.
% 3. On compare 'c' à 'max'. Si 'c' est plus grand, 'max' prend la valeur de 'c'.
% 4. À la fin, 'max' contient la plus grande des trois valeurs.
\begin{tikzpicture}[node distance=1.5cm and 2cm]
    \node (start) [startstop] {Début};
    \node (in1)   [io, below=of start] {Saisir a, b, c};
    \node (proc1) [process, below=of in1] {max $\leftarrow$ a};
    \node (dec1)  [decision, below=of proc1, yshift=-0.5cm] {b $>$ max ?};
    \node (proc2) [process, right=of dec1] {max $\leftarrow$ b};
    
    % Point de jonction après la première comparaison
    \coordinate (join1) at ([below=1.2cm]dec1.south);

    \node (dec2)  [decision, at=(join1), yshift=-1.2cm] {c $>$ max ?};
    \node (proc3) [process, right=of dec2] {max $\leftarrow$ c};

    % Point de jonction après la deuxième comparaison
    \coordinate (join2) at ([below=1.2cm]dec2.south);
    
    \node (out1)  [io, at=(join2), yshift=-1.2cm] {Afficher max};
    \node (stop)  [startstop, below=of out1] {Fin};

    % Flèches du flux principal
    \draw [arrow] (start) -- (in1);
    \draw [arrow] (in1) -- (proc1);
    \draw [arrow] (proc1) -- (dec1);

    % Branche et jonction pour la comparaison de 'b'
    \draw [arrow] (dec1.east) -- node[above] {Oui} (proc2);
    \draw [arrow] (dec1.south) -- node[left] {Non} (join1);
    \draw [arrow] (proc2) |- (join1);

    \draw [arrow] (join1) -- (dec2);

    % Branche et jonction pour la comparaison de 'c'
    \draw [arrow] (dec2.east) -- node[above] {Oui} (proc3);
    \draw [arrow] (dec2.south) -- node[left] {Non} (join2);
    \draw [arrow] (proc3) |- (join2);

    \draw[arrow] (join2) -- (out1);
    \draw[arrow] (out1) -- (stop);
\end{tikzpicture}
\end{center}

\end{document}